\section{Self-Improvement Methods}
\label{sec:self-improvement}

The capacity for self-improvement is the most fundamental form of self-evolution in AI systems. In this section, we provide a comprehensive taxonomy and analysis of methods through which large language models (LLMs) and language agents improve their own performance without explicit human intervention. We organize these methods along two principal axes: \textbf{inference-time self-improvement}, where improvement occurs within a single forward pass or a bounded number of interaction rounds without modifying model weights; and \textbf{training-time self-improvement}, where the model's parameters are iteratively refined using self-generated data.

This distinction is consequential. Inference-time methods offer immediate, training-free improvement at the cost of additional compute per query, while training-time methods produce persistent improvements that amortize over many queries but require computational investment in fine-tuning. As we demonstrate below, the most impactful recent systems increasingly combine both paradigms.

\subsection{Inference-Time Self-Improvement}
\label{sec:inference-time}

Inference-time self-improvement methods enable a model to produce better outputs than its initial generation through iterative refinement, reflection, or debugging---all without modifying the model's parameters. These methods treat the LLM as a fixed function and improve outputs through structured interaction protocols.

\subsubsection{Self-Refine: Iterative Refinement with Self-Feedback}
\label{sec:self-refine}

Self-Refine~\citep{madaan2023selfrefine} introduces a minimal yet powerful framework: a single LLM iteratively generates, critiques, and refines its own output. The method requires no external feedback, no reinforcement learning, and no weight updates.

\paragraph{Method.}
The Self-Refine loop consists of three stages applied iteratively:

\begin{enumerate}
    \item \textbf{GENERATE}: The model produces an initial output $y_0 = \text{LLM}(x, \text{instruction})$ given input $x$.
    \item \textbf{CRITIQUE}: The model analyzes its own output and produces structured feedback: $\text{fb}_t = \text{LLM}(y_t, \text{rubric})$, where the rubric defines task-specific evaluation criteria.
    \item \textbf{REFINE}: The model improves the output based on the critique: $y_{t+1} = \text{LLM}(y_t, \text{fb}_t, \text{instruction})$.
\end{enumerate}

The loop terminates when the critique indicates no further improvements are needed or a maximum iteration count $T$ is reached. A critical design choice is that the \emph{same model} serves all three roles, making Self-Refine applicable to any sufficiently capable LLM without additional training.

\paragraph{Results.}
Self-Refine achieves approximately 20\% average improvement across seven diverse tasks~\citep{madaan2023selfrefine}. On code optimization, it achieves a 28\% improvement; on math reasoning, accuracy rises from 56\% to 67\% (+11 percentage points); on sentiment reversal, accuracy increases from 74\% to 86\% (+12pp); on dialogue response, win rate rises from 52\% to 68\% (+16pp); on constrained generation, compliance increases from 65\% to 78\% (+13pp); on acrostic poetry, quality improves from 3.2 to 3.8 (+0.6); and on code readability, score improves from 3.5 to 4.1 (+0.6).

\paragraph{Limitations.}
Self-Refine is bounded by the quality of the initial generation: if the model's first attempt is fundamentally flawed, self-critique may fail to identify the core issue. The method can also converge to local optima, where iterative refinement makes superficial improvements without addressing deeper structural problems. Each iteration increases compute cost linearly, and the method may ``refine away'' creative or diverse outputs.

\subsubsection{Reflexion: Verbal Reinforcement Learning}
\label{sec:reflexion}

Reflexion~\citep{shinn2023reflexion} extends the self-improvement paradigm by introducing persistent linguistic memory that accumulates across episodes. Rather than refining a single output, Reflexion learns from failed attempts on a task to improve performance on subsequent attempts.

\paragraph{Method.}
Reflexion employs a three-model architecture:

\begin{itemize}
    \item \textbf{Actor} ($M_a$): Generates actions or outputs given the task and memory context.
    \item \textbf{Evaluator} ($M_e$): Scores the actor's output, producing a binary reward $r_e \in \{0, 1\}$ or a scalar reward.
    \item \textbf{Self-Reflection} ($M_r$): Generates verbal feedback explaining what went wrong and how to improve.
\end{itemize}

The Reflexion loop operates as follows. For each episode $e = 1, 2, \ldots, E$:

\begin{align}
    a_e &= M_a(\text{task}, M_{1:e-1}) \label{eq:reflexion-actor} \\
    r_e &= M_e(a_e, \text{task}) \label{eq:reflexion-eval} \\
    f_e &= M_r(a_e, r_e, \text{task}) \quad \text{if } r_e \text{ indicates failure} \label{eq:reflexion-reflect} \\
    M_e &= M_{e-1} \cup \{f_e\} \label{eq:reflexion-memory}
\end{align}

The key innovation is the memory buffer $M$, which stores linguistic feedback across episodes. Each entry describes what the agent tried, why it failed, and what to do differently. This acts as a ``verbal gradient''---natural language guidance that replaces the role of parameter gradients in traditional reinforcement learning.

\paragraph{Results.}
Reflexion achieves 91\% pass@1 on HumanEval using GPT-3.5, surpassing GPT-4's 80\% in the zero-shot setting~\citep{shinn2023reflexion}. On AlfWorld (interactive environments), Reflexion achieves 97\% success versus 75\% for the ReAct baseline. On MBPP (code generation), significant improvements are observed. The method demonstrates clear scaling behavior: more reflection episodes yield better performance, with specific feedback outperforming generic feedback. Cross-task transfer is observed: feedback from similar tasks helps.

\paragraph{Limitations.}
Reflexion requires a scalar or binary evaluation signal, which may not be available for all tasks. The memory buffer grows unbounded, eventually exceeding the model's context window. The quality of self-reflection depends on the underlying LLM's capability, and there is no theoretical framework for convergence guarantees.

\subsubsection{Self-Debugging: Execution-Guided Code Repair}
\label{sec:self-debug}

Self-Debugging~\citep{chen2023selfdebug} specializes self-improvement to the code domain by incorporating execution feedback as an external verification signal. The key innovation is ``rubber duck debugging''---the model explains its own code line by line, often discovering errors during the explanation process.

\paragraph{Method.}
The Self-Debug loop operates as:

\begin{enumerate}
    \item \textbf{GENERATE}: $p_0 = \text{LLM}(x, \text{prompt})$---the model generates an initial program.
    \item \textbf{EXECUTE}: $f_t = \text{Execute}(p_t, \text{test\_cases})$---the program is executed against test cases.
    \item \textbf{EXPLAIN}: $e_t = \text{LLM}(p_t, \text{``explain this code''})$---the model explains its own code (rubber duck debugging).
    \item \textbf{DEBUG}: $p_{t+1} = \text{LLM}(p_t, f_t, e_t, \text{debug\_prompt})$---the model repairs the program.
\end{enumerate}

Three types of feedback can be provided: (1)~Simple Feedback (pass/fail signal), (2)~Unit Test Feedback (input/output pairs from failed tests), and (3)~Code Explanation (the model's own explanation). Crucially, the explanation-based approach requires \emph{no external execution environment}.

\paragraph{Results.}
On Spider (text-to-SQL), Self-Debug with code explanation achieves up to 9\% improvement on the hardest difficulty level. On TransCoder (C++ to Python) and MBPP (Python generation), up to 12\% accuracy improvement. Self-Debug is far more sample-efficient than best-of-$K$ sampling, matching baselines that generate 10$\times$ more candidates.

\paragraph{Significance.}
Self-Debug bridges inference-time and execution-time improvement by leveraging a deterministic external evaluator (the code interpreter). This paradigm---generate, execute, diagnose, repair---recurs throughout the evolutionary code systems analyzed in Section~\ref{sec:evolutionary-code}.

\subsubsection{Agent Symbolic Learning: Textual Backpropagation}
\label{sec:symbolic-learning}

Agent Symbolic Learning~\citep{zhou2024symbolic} represents a significant conceptual advance: it treats an LLM-based agent as a \emph{symbolic network} whose ``weights'' are natural language prompts, tool descriptions, and pipeline compositions, and then applies an analog of backpropagation and gradient descent entirely in text space.

\paragraph{Method.}
An agent is parameterized as a computational graph with three types of learnable nodes: \textbf{Prompt nodes} ($\mathcal{P}$) for system prompts and instructions, \textbf{Tool nodes} ($\mathcal{T}$) for function definitions, and \textbf{Pipeline nodes} ($\mathcal{I}, \mathcal{O}$) for input/output connections.

The framework operates in four stages: (1)~\textbf{Forward Pass}: execute the agent on a task to produce trajectory $\tau$; (2)~\textbf{Language Loss}: $\mathcal{L}_{\text{lang}} = \text{LLM}(\mathcal{P}_{\text{loss}}(\tau))$; (3)~\textbf{Language Gradient Backpropagation}: for each layer $n$ from output to input:
\begin{equation}
    \nabla_{\text{lang}}^n = \text{LLM}(\mathcal{P}_{\text{gradient}}(\nabla_{\text{lang}}^{n+1}, \mathcal{I}_n, \mathcal{O}_n, \mathcal{P}_n, \mathcal{T}_n, \mathcal{L}_{\text{lang}}))
\end{equation}
(4)~\textbf{Weight Update}: PromptOptimizer updates $\Delta\mathcal{P}$, ToolOptimizer updates $\Delta\mathcal{T}$, and PipelineOptimizer updates $\Delta\mathcal{I}, \Delta\mathcal{O}$.

\paragraph{Results.}
On HotPotQA, Agent Symbolic Learning achieves F1 = 39.1 and EM = 25.6, outperforming both vanilla agents (F1 = 36.2) and DSPy (F1 = 37.4). On MATH with GPT-4 backbone, accuracy reaches 60.7\% versus 56.0\% for baseline agents and 48.4\% for DSPy. On HumanEval, 73.2\% pass@1 versus 68.3\% baseline. For software development (1--4 executability score), 3.8 versus 2.4 for baseline agents and 1.6 for GPTs. Creative writing: 7.4 versus 6.8 for Tree of Thought.

\paragraph{Significance.}
Agent Symbolic Learning is the first framework to apply structured optimization---analogous to neural network training---to the full architecture of an LLM agent. It demonstrates that agent ``weights'' (prompts, tools, pipelines) can be systematically improved through a principled learning procedure, establishing a foundation for the more radical self-modifying systems discussed in Section~\ref{sec:agent-evolution}.

\subsection{Training-Time Self-Improvement}
\label{sec:training-time}

Training-time self-improvement methods modify the model's parameters using self-generated data, producing persistent improvements. These methods follow a common generate-filter-finetune paradigm but differ in how they generate data, filter for quality, and perform the update.

\subsubsection{STaR: Self-Taught Reasoning}
\label{sec:star}

STaR (Self-Taught Reasoner)~\citep{zelikman2022star} pioneered the generate-filter-finetune paradigm for LLM self-improvement. The key insight is that a model can bootstrap its own reasoning capability by generating reasoning chains (rationales), filtering for correct answers, and fine-tuning on successful chains.

\paragraph{Method.}
Given a dataset $D = \{(x_i, y_i)\}$ and demonstrations $F$, STaR iterates: (1)~\textbf{GENERATE}: $(r_i, \hat{y}_i) = \arg\max_{r,y} p_\theta(r, y | x_i, F)$; (2)~\textbf{FILTER}: $D_{\text{correct}} = \{(x_i, r_i) : \hat{y}_i = y_i\}$; (3)~\textbf{RATIONALIZE}: For incorrect examples, provide the answer as a hint and regenerate $r_i'$; add successful rationalizations to $D_{\text{correct}}$; (4)~\textbf{FINETUNE}: $\mathcal{L} = -\sum_{(x,r) \in D_{\text{correct}}} \log p_\theta(r | x, F)$.

\paragraph{Rationalization.}
The rationalization step is STaR's key innovation. Rather than discarding failed examples, STaR tells the model the correct answer and asks it to ``work backward'' to construct a valid reasoning chain, significantly increasing effective training set size.

\paragraph{Results.}
On CommonsenseQA, STaR achieves 72.5\% accuracy using a 540M parameter model, comparable to a fine-tuned 175B model (30$\times$ larger). Adding rationalization contributes $\sim$5\% additional improvement, with larger gains on harder problems. The method demonstrates that self-improvement can close a 300$\times$ scale gap.

\subsubsection{SPIN: Self-Play Fine-Tuning}
\label{sec:spin}

SPIN~\citep{chen2024spin} frames LLM self-improvement as a two-player game between two versions of the same model. The training objective is:
\begin{equation}
    L_{\text{SPIN}}(\theta) = \mathbb{E}_{x \sim D}\left[\mathbb{E}_{y \sim p_{\text{data}}}\left[\log \sigma\left(\lambda(h_\theta(x,y) - h_\theta(x,\tilde{y}))\right)\right]\right]
    \label{eq:spin}
\end{equation}
where $\tilde{y} \sim p_{\theta_t}(\cdot|x)$ is the previous model's response and $y \sim p_{\text{data}}$ is the human response.

A notable theoretical contribution is the proof that the global optimum satisfies $p_{\theta^*} = p_{\text{data}}$. Convergence occurs in approximately 3 iterations. SPIN, using only SFT data without additional preference data, outperforms DPO trained with GPT-4 preference data on multiple benchmarks.

\subsubsection{ReST-EM: Reinforced Self-Training}
\label{sec:rest-em}

ReST-EM~\citep{gulcehre2023rest} formalizes generate-filter-finetune as an EM procedure. The E-step generates $N$ candidates using the current policy; the M-step scores them using a reward model, filters by threshold, and fine-tunes on filtered data. On machine translation, ReST achieves 1.5--3.0 BLEU point improvements and outperforms baselines in human evaluations.

\subsubsection{Constitutional AI: Self-Alignment through Principles}
\label{sec:constitutional-ai}

Constitutional AI~\citep{bai2022constitutional} introduces self-improvement for alignment: the model generates responses to harmful prompts, critiques itself against constitutional principles, revises responses, and uses revised outputs for fine-tuning. A second RL phase uses model-generated preference data to train a reward model. This demonstrates that self-improvement need not be limited to task performance---it can also improve safety and alignment.

\subsection{Hybrid Methods: Bridging Inference and Training}
\label{sec:hybrid-methods}

The most promising direction integrates inference-time and training-time self-improvement. Two recent systems exemplify this hybrid approach.

\subsubsection{RISE: Recursive Introspection via RL}
\label{sec:rise}

RISE~\citep{qu2024rise} formulates self-improvement as a multi-turn MDP where the model learns when and how to self-correct. The policy is trained to maximize:
\begin{equation}
    \max_\theta \mathbb{E}_{\pi_\theta}[R] - \beta \, \text{KL}(\pi_\theta \| \pi_{\text{ref}})
    \label{eq:rise-objective}
\end{equation}

On GSM8K, RISE improves Llama2-7B from $\sim$14\% to $\sim$24\% (+10pp) and Mistral-7B from $\sim$38\% to $\sim$46\% (+8pp). Critically, RL fine-tuning significantly outperforms Self-Refine-style prompting, demonstrating that self-correction is more effective when \emph{learned} than when \emph{instructed}. The model exhibits emergent difficulty adaptation: it allocates more self-correction budget to harder problems.

RISE demonstrates a key principle: self-improvement capabilities can be \emph{trained into} models rather than merely prompted. This insight underpins many advanced systems in subsequent sections.

\subsubsection{RAGEN: Multi-Turn Agent RL and the Echo Trap}
\label{sec:ragen}

RAGEN~\citep{wang2025ragen} extends RL-based self-improvement to the full multi-turn agent setting with StarPO (State-Thinking-Actions-Reward Policy Optimization):
\begin{equation}
    \max_\theta \mathbb{E}\left[\sum_t \gamma^t \cdot r_t \cdot \log \pi_\theta(a_t | s_t, t_t)\right]
    \label{eq:starpo}
\end{equation}

RAGEN's most important contribution is identifying the ``Echo Trap'': during multi-turn RL training, agents learn to \emph{echo} previous outputs rather than genuinely reasoning. Detection: $\text{cos\_sim}(t_t, t_{t-1}) > \theta_{\text{echo}}$. The proposed solution, StarPO-S, filters echo-dominated trajectories.

The Echo Trap reveals a fundamental challenge: without careful monitoring, agents can ``fake'' improvement by exploiting patterns in their own outputs. This has direct implications for the evaluation and safety of self-evolving systems.

\subsubsection{ReVeal: Co-Evolution of Generation and Verification}
\label{sec:reveal}

ReVeal~\citep{jin2025reveal} (ICLR 2026, Microsoft Research) explicitly co-evolves code generation and self-verification. The TAPO framework provides per-turn credit assignment:
\begin{equation}
\tau = \{(c_1, v_1, r_1), \ldots, (c_T, v_T, r_T)\}
\end{equation}
where $c_t$ is code, $v_t$ is self-verification, and $r_t$ is the turn-level reward. The loss jointly optimizes: $\mathcal{L} = \mathcal{L}_{\text{generation}} + \lambda \times \mathcal{L}_{\text{verification}}$.

Strikingly, ReVeal trained with only 3 turns can self-evolve for 20+ turns at inference---demonstrating that learned self-correction generalizes far beyond the training horizon. This ``turn extrapolation'' suggests the RL training process learns a general self-improvement policy.

\subsection{Autonomous Scientific Discovery}
\label{sec:ai-scientist}

AI Scientist~\citep{lu2024aiscientist} (Sakana AI) represents the most ambitious self-improvement application: fully automated scientific research. The system performs the entire lifecycle---idea generation, experimental design, code execution, manuscript writing, and automated peer review---without human intervention.

AI Scientist-v2 achieved the first entirely AI-generated paper accepted through peer review at a workshop level (score 6.33, top 45\%). However, independent evaluation shows the LLM-based reviewer rejected 9/10 papers including 4 human-accepted ones, and the system struggles to generalize beyond narrow training domains. This demonstrates the evaluation bottleneck that limits all self-improvement methods.

\subsection{Comparative Analysis}
\label{sec:comparison}

Table~\ref{tab:self-improvement-comparison} provides a systematic comparison across key dimensions.

\begin{table}[t]
\centering
\caption{Comparison of self-improvement methods. ``Feedback Source'' indicates where the improvement signal originates.}
\label{tab:self-improvement-comparison}
\small
\begin{tabular}{@{}lcccc@{}}
\toprule
\textbf{Method} & \textbf{Timing} & \textbf{Weight Update} & \textbf{Feedback Source} & \textbf{Key Innovation} \\
\midrule
Self-Refine & Inference & No & Self & Single-model refine loop \\
Reflexion & Inference & No & Self + Evaluator & Verbal memory buffer \\
Self-Debug & Inference & No & Execution & Rubber duck debugging \\
Symbolic Learning & Inference & No (text) & Self & Textual backprop \\
RISE & Both & RL & Self + Reward & Learned self-correction \\
RAGEN & Training & RL & Self + Reward & Echo Trap detection \\
ReVeal & Training & RL & Self + Verification & Gen-verify co-evolution \\
\midrule
STaR & Training & SFT & Human labels & Rationalization \\
SPIN & Training & SFT & Human data & Self-play convergence \\
ReST-EM & Training & SFT & Reward model & EM framework \\
Constitutional AI & Training & SFT+RL & Principles & Self-alignment \\
\bottomrule
\end{tabular}
\end{table}

\paragraph{The Generate-Filter-Finetune Paradigm.}
A striking pattern: nearly all training-time methods follow the same loop---generate candidates, filter for quality, finetune on filtered data. This mirrors the evolutionary generate-select-mutate loop explored in Section~\ref{sec:evolutionary-code}. The key variable is the filtering mechanism: STaR uses answer correctness, SPIN uses distinguishability from human data, ReST-EM uses a reward model, and Constitutional AI uses principled self-critique.

\paragraph{Scaling Relationships.}
\textbf{Scale compression}: STaR's 540M matches 175B, suggesting self-improvement is most impactful for smaller models. \textbf{Iteration convergence}: Most methods converge in 2--5 iterations (SPIN $\sim$3, ReST-EM 2--3, STaR 3--5). \textbf{Compute-performance tradeoff}: Inference-time methods scale with per-query compute; training-time methods scale with upfront compute.

\paragraph{Connections to the Five Evolution Loops.}
Each method instantiates one or more taxonomy loops (Section~\ref{sec:taxonomy}): \textbf{Reflection Loop}---Self-Refine, Reflexion, Agent Symbolic Learning, Constitutional AI; \textbf{Evaluator Loop}---Self-Debug, RISE, ReST-EM; \textbf{Population Loop}---SPIN's self-play; \textbf{Search Loop}---Agent Symbolic Learning's textual backpropagation.

\subsection{Benchmark Data Compilation}
\label{sec:benchmark-compilation}

Table~\ref{tab:benchmark_compilation} consolidates quantitative results.

\begin{table}[t]
\centering
\small
\caption{Benchmark results compilation. All results from original papers.}
\label{tab:benchmark_compilation}
\begin{tabular}{llcc}
\toprule
\textbf{Benchmark} & \textbf{Method} & \textbf{Base} & \textbf{Improved} \\
\midrule
\multirow{3}{*}{HumanEval (Pass@1)} & GPT-4 (zero-shot) & 80.1\% & --- \\
& Reflexion (GPT-3.5) & 67.0\% & \textbf{91.0\%} \\
& Agent Symbolic Learning & 68.3\% & 73.2\% \\
\midrule
\multirow{2}{*}{CommonsenseQA} & STaR (540M) & 53.3\% & \textbf{72.5\%} \\
& Fine-tuned 175B & --- & ~73\% \\
\midrule
\multirow{2}{*}{GSM8K} & RISE (Mistral-7B) & ~38\% & \textbf{~46\%} \\
& RISE (Llama2-7B) & ~14\% & ~24\% \\
\midrule
\multirow{2}{*}{HotPotQA (F1)} & Agent Symbolic Learning & 36.2 & \textbf{39.1} \\
& DSPy & 37.4 & --- \\
\midrule
\multirow{2}{*}{MATH (Accuracy)} & Agent Symbolic Learning & 56.0\% & \textbf{60.7\%} \\
& DSPy & 48.4\% & --- \\
\midrule
\multirow{2}{*}{AlfWorld (Success)} & Reflexion & 75\% & \textbf{97\%} \\
& Base ReAct & 75\% & --- \\
\midrule
Code Optimization & Self-Refine & Baseline & \textbf{+28\%} \\
Sentiment Reversal & Self-Refine & 74\% & \textbf{86\%} \\
Constrained Gen & Self-Refine & 65\% & \textbf{78\%} \\
\bottomrule
\end{tabular}
\end{table}

The most striking finding is \textbf{size-equivalence}: self-improvement enables small models to match or exceed large models (STaR's 540M matching 175B; Reflexion's GPT-3.5 surpassing GPT-4), suggesting that self-improvement can substitute for scale.

\subsection{Extended Analysis of Self-Improvement Dynamics}
\label{sec:extended-analysis}

\subsubsection{The Self-Correction Paradox}

A counter-intuitive finding is that naive self-correction---simply asking a model to ``check your answer''---often \emph{degrades} performance. The methods that overcome this share a common feature: they introduce \textbf{new information}. Reflexion introduces evaluator feedback; Self-Debug introduces execution traces; RISE introduces RL-learned correction behaviors; Agent Symbolic Learning introduces language-space gradients. The implication: self-improvement requires external grounding. Systems relying solely on self-reference without external feedback are likely to oscillate or degrade.

\subsubsection{Scaling Laws for Self-Improvement}

\textbf{Model size scaling}: Self-improvement is most effective for smaller models (STaR's 300$\times$ scale compression). Very large models produce near-optimal outputs, leaving less room for improvement. \textbf{Iteration scaling}: Convergence in 2--5 iterations suggests naive self-training has limited capacity for sustained improvement. \textbf{Domain scaling}: Strongest results in structured domains (code, math, reasoning with ground truth); weaker in open-ended tasks.

\subsubsection{The Inference-Training Boundary}

RISE and RAGEN blur the traditional boundary. RISE learns self-correction at inference time through RL training; RAGEN trains multi-turn agents that improve within episodes through learned self-reflection. The key insight: self-correction is a \emph{learned skill} that can be trained. RAGEN's Echo Trap adds a caveat: learned self-correction can be gamed when training rewards apparent improvement.

\subsection{Cross-Cutting Observations}
\label{sec:cross-cutting}

\begin{enumerate}
    \item \textbf{The generate-filter-train pattern is universal.} Every training-time method follows this structure, differing only in filtering mechanism and training objective.
    \item \textbf{Inference-time and training-time methods are complementary.} Use inference-time methods for immediate gains; periodically distill into parameters through training-time methods.
    \item \textbf{The evaluation bottleneck persists.} Quality of self-improvement is bounded by quality of evaluation. The Echo Trap demonstrates that naive self-evaluation can be actively harmful.
    \item \textbf{Convergence is fast (2--3 iterations).} Sustained improvement requires diverse training signals or increasingly challenging evaluation criteria.
    \item \textbf{New information is required for effective self-correction.} Every successful method introduces new information---from execution, evaluation, or learned policies.
    \item \textbf{Self-correction is a learnable skill.} RISE and ReVeal demonstrate that the ability to self-correct can be trained via RL, producing improvement that prompting alone cannot achieve.
    \item \textbf{Turn extrapolation enables inference-time scaling.} ReVeal's 3-turn training enabling 20+ turn self-evolution suggests learned improvement policies generalize beyond their training horizon.
\end{enumerate}

\subsection{Transitions to Subsequent Sections}
\label{sec:transitions}

The methods in this section primarily improve model \emph{outputs}---responses, code, reasoning chains. The improvement mechanisms operate within fixed architectures. In Section~\ref{sec:evolutionary-code}, we examine systems that improve the \emph{programs} that AI executes, discovering solutions surpassing decades of human effort. In Section~\ref{sec:agent-evolution}, we examine systems that improve the \emph{agents} themselves---their architectures, workflows, and learning mechanisms---representing the most radical form of self-improvement: agents that modify their own code and evolve their own architectures.
